%═══════════════════════════════════════════════════════════════════════════════
\section{Application: Stocking Shirts for the World Cup Final}\label{sec:market}
%═══════════════════════════════════════════════════════════════════════════════

Every example so far has stipulated its distributions. This section applies
the procedure of Section~\ref{sec:recipe} to real data, taking the belief from
Polymarket prices for the 2026 FIFA World Cup final between Spain and
Argentina \citep{Marketlens2026}. The decision-maker is a shopkeeper near the
stadium who orders replica shirts before the final and cannot reorder
afterwards. Shirts of the winning side sell and the rest are marked down. The
shopkeeper does not bet on the match and uses the prices only as a forecast,
so the cost of an order depends on how well it matches demand, never on the
prices themselves.

The section reaches one finding. Under the market's forecast the cheapest
order is mostly Spain shirts, yet no increasing reshaping of the likelihood
ratio leads to that order, across two extensive searches, while the safe
profile does. The reason is where each places its evidence. Protecting a
Spain-heavy order needs large scores on the Argentina outcomes, where that
order is ruinous. A reshaped likelihood ratio puts its largest score on a
Spain win within ninety minutes, where the market is most confident and the order
costs almost nothing. This is Result 2 in real data: the profile that best
reflects the belief leads away from the action the belief recommends.

\subsection{Two Contracts, Four Outcomes}\label{sec:mkt_data}

Two contracts were traded on the final, each settling on a different
question. A three-way result market pays out on the score after ninety
minutes, so a match still level then settles as a draw. A binary trophy
market pays out on whichever side eventually lifts the cup. Together they
tell the shopkeeper which side wins and whether the match is decided within
ninety minutes, and both affect how many shirts sell.

Figure~\ref{fig:market_odds} shows the two on one timeline. Prices are
recorded every minute of clock time, so breaks in play appear as gaps, which
is why the goal in the $106$th minute of play falls more than two hours after
kick-off. The two series part at the end of regulation. With the match level
the draw contract settles and the result market stops moving, while the
trophy market stays live and Spain's extra-time goal lifts their price from
about $0.60$ to over $0.90$. These are traded prices, not e-values.

\begin{figure}[H]
\centering
\includegraphics[width=0.92\textwidth]{polymarket_match_odds.pdf}
\caption{Two Polymarket contracts on the same match, from kick-off to settlement. The dashed line marks the end of regulation, the dotted line the goal.}\label{fig:market_odds}
\end{figure}

At kick-off the result market quoted
\[
  q=(0.4231,\;0.3238,\;0.2531)
\]
for a Spain win, a draw and an Argentina win, and the trophy market quoted
$o=(0.5863,\,0.4137)$ for Spain and Argentina. Each number is the midpoint of
the best bid and ask. Midpoints of separately traded contracts need not sum to
one, so each set is rescaled to sum to one, here and in
Figure~\ref{fig:market_odds}. The data also record how both contracts
settled: the draw contract paid out, and the trophy went to Spain
\citep{Marketlens2026}.

Read together, the two contracts describe four outcomes,
\[
  \mathcal Z=\{s,\;s^{\mathrm e},\;a,\;a^{\mathrm e}\},
\]
where $s$ is a Spain win within ninety minutes, $s^{\mathrm e}$ a Spain win
after extra time or penalties, and $a$ and $a^{\mathrm e}$ the same for
Argentina. Four outcomes need three numbers, and the two contracts supply
exactly three. The result market gives the two wins within ninety minutes
directly, and subtracting these from the trophy prices leaves the probability
of each side winning later:
\begin{equation}
  \widehat P_1(s)=q_1,\quad
  \widehat P_1(s^{\mathrm e})=o_1-q_1,\quad
  \widehat P_1(a)=q_3,\quad
  \widehat P_1(a^{\mathrm e})=o_2-q_3.
  \label{eq:market_ident}
\end{equation}
A match level after ninety minutes must eventually be won by someone, so the
two late wins should together account for the draw price, and they do:
$0.1632+0.1606=0.3238=q_2$. This holds by construction, since both sets of
prices were rescaled to sum to one, so it checks the arithmetic rather than
the market. The contracts cannot tell extra time from penalties, and nothing
below depends on that distinction.

These prices give the belief only. The model, the set of distributions the
guarantee must hold under, is not read from the prices at all. It is built
separately in the next subsection.

\subsection{The Model and the Belief}\label{sec:mkt_dists}

The procedure needs a belief and a model, and they do different jobs. The
belief is the shopkeeper's best guess at what will happen, and it decides
which order looks cheapest. The model is the set of distributions the
guarantee must hold up against, and it decides how strong a guarantee any
profile can give. Both must be fixed before kick-off, when the order is
placed.

The belief is the market's view at kick-off. Applying \eqref{eq:market_ident}
to the opening quotes gives
\[
  \widehat P_1=(0.4231,\;0.1632,\;0.2531,\;0.1606).
\]
Later prices cannot replace it. They absorb the match as it is played, and by
the end nearly all their probability sits on the outcome that happened, so
they report the result rather than forecast it.

The model needs more explanation. The guarantee is meant to hold even if the
market is wrong about this match, so the model cannot simply be the market's
forecast, or the guarantee would be only as good as the market. Instead it
collects a few distributions built from information available before the two
finalists were known, and the guarantee must hold under all of them.
These are not meant to be better forecasts than the market, which is better
informed. They are the alternatives the shopkeeper wants to be protected
against.

Each member splits the probability between the two sides, and then splits
each side between winning within ninety minutes and winning later. Three
splits between the sides are used: Polymarket's tournament-winner prices
early in the group stage, $0.138$ for Spain and $0.116$ for Argentina
\citep{PolymarketWinner2026}, rescaled to sum to one as
$(0.5433,0.4567)$, an even split, and $(0.45,0.55)$, leaning to Argentina. The share of matches decided later
comes from history. Of the twenty-one World Cup finals played as a single
deciding match between 1930 and 2022, eight went beyond ninety minutes
\citep{WCfinalsList}, a rate of $8/21=0.381$. The 1950 tournament is excluded
because it ended with a final group stage rather than a final. For the first
split this gives
\[
\begin{aligned}
  P&=\bigl(0.5433\,(0.619),\;0.5433\,(0.381),\;
             0.4567\,(0.619),\;0.4567\,(0.381)\bigr)\\
   &=(0.3363,\;0.2070,\;0.2827,\;0.1740),
\end{aligned}
\]
and the other two follow in the same way, so
\[
\begin{aligned}
  \mathcal P=\bigl\{&(0.3363,\;0.2070,\;0.2827,\;0.1740),\\
                    &(0.3095,\;0.1905,\;0.3095,\;0.1905),\\
                    &(0.2786,\;0.1714,\;0.3405,\;0.2095)\bigr\}.
\end{aligned}
\]

The model is deliberately more cautious than the belief. The market puts
$0.4231$ on a Spain win within ninety minutes, while no member of the model
goes above $0.3363$. By Theorem~\ref{thm:cert}, only one member matters for a
given order, the one under which that order looks worst, and
Section~\ref{sec:mkt_shape} checks that the findings survive changes to the
historical rate.

\subsection{The Two Criteria Disagree}\label{sec:mkt_decision}

The shopkeeper has room for one hundred shirts and must split them between the
two sides. Three orders are considered: mostly Spain ($d_1$, ninety Spain
shirts), an even split ($d_2$, fifty) and mostly Argentina ($d_3$, twenty).
Each shirt short of demand costs the lost margin, three units, and each shirt
left unsold costs the markdown, one unit.

Demand depends on the outcome. Only the winning side's shirts sell well, and a
win after extra time is a bigger occasion, so it sells more. Demand for the
winner's shirts is taken as eighty after a win within ninety minutes and one
hundred after a later win, with ten of the losing side either way. The cost of
each order then follows by arithmetic. In tens of units, for the outcomes
$(s,s^{\mathrm e},a,a^{\mathrm e})$,
\[
 L(d_1,\cdot)=(1,3,29,35),\quad
 L(d_2,\cdot)=(13,19,13,19),\quad
 L(d_3,\cdot)=(25,31,1,7).
\]

For example, ordering mostly Spain when Spain win within ninety minutes leaves
ten Spain shirts unsold, a cost of ten units, or $1$ in the table. Ordering
mostly Spain is cheap when Spain win and ruinous when Argentina do. The even
split is never cheap and never ruinous, and its cost depends only on when the
match is decided, not on who wins. No price appears anywhere in the table: the
market says what is likely, and the shirts say what it costs to be wrong.

Step one compares the orders in the two ways of Section~\ref{sec:recipe}, by
worst-case expected loss across the model and by expected loss under the
belief,
\[
\begin{aligned}
  \max_{P\in\mathcal P}\mathbb E_{P}[L(\cdot,Z)]&=(18.0000,\;15.2857,\;16.3251),\\
  \mathbb E_{\widehat P_1}[L(\cdot,Z)]&=(13.8738,\;14.9430,\;17.0142).
\end{aligned}
\]
The two disagree. The protective action is the even split $d_2$, the order
that cannot go badly wrong. The belief recommends mostly Spain, $d_1$, because
the market makes Spain the likelier winner. This is the case the thesis is
about.

Step two takes $d_1$ as the target, the choice of a belief-trusting
shopkeeper. The worst member of the model for $d_1$ is the one leaning to
Argentina,
\[
  P^\star=(0.2786,\;0.1714,\;0.3405,\;0.2095),
\]
under which
\[
  \mu_{d_1}=1\,(0.2786)+3\,(0.1714)+29\,(0.3405)+35\,(0.2095)=18.0000,
\]
so the safe profile of Theorem~\ref{thm:cert} divides the losses of $d_1$ by
$18$,
\begin{equation}
  \mathcal E^{\mathrm{safe}}_{d_1}
  =\frac{(1,3,29,35)}{18}
  =(0.0556,\;0.1667,\;1.6111,\;1.9444).
  \label{eq:market_safe}
\end{equation}
It puts most of its e-value on the two Argentina outcomes, where $d_1$ is most
exposed. Its certificate is $18$, and since $d_1$ is the action the belief
recommends, Corollary~\ref{cor:select} guarantees that it leads to $d_1$.
Table~\ref{tab:market} compares it with the profiles an evidence-first
statistician would report. The likelihood ratio there is taken against
$P^\star$ and scaled down by a factor of $1.051$, so that it is valid under
every member of the model.

\begin{table}[htbp]
\centering
\caption{The shopkeeper's order. Certificates $C_j$ of each order under each profile, the order the rule selects, and that order's expected loss under the kick-off belief.}\label{tab:market}
\begin{tabular}{@{}l r r r c r@{}}\toprule
Profile & $C_1$ & $C_2$ & $C_3$ & Action & $\mathbb E_{\widehat P_1}[L]$ \\\midrule
Likelihood ratio          & 48.002 & 26.058 & 34.233 & $d_2$ & 14.943 \\
Flat                      & 35.000 & 19.000 & 31.000 & $d_2$ & 14.943 \\
Safe for $d_2$            & 34.099 & 15.286 & 29.396 & $d_2$ & 14.943 \\
Safe for $d_1$            & 18.000 & 234.000 & 450.000 & $d_1$ & 13.874 \\\bottomrule
\end{tabular}
\end{table}

Only the safe profile for $d_1$ leads to $d_1$. The likelihood ratio and the
flat profile both leave the shopkeeper with the even split, and so does the
safe profile for $d_2$. This is the failure Section~\ref{sec:result2} warned
of: the likelihood ratio does not lead to the action the belief itself
recommends. Measured under the belief, that costs the shopkeeper about $1.07$
in expected loss, $14.943$ against $13.874$. The prediction-set rule of
Section~\ref{sec:intro_sets} does no better: at every coverage level from
$50\%$ to $95\%$ it also chooses the even split.

\subsection{No Increasing Shape Reaches the Target}\label{sec:mkt_shape}

Only the safe profile led to the order the belief recommends. A natural
objection is that the standard construction was not tried hard enough, and
that some other shape function would have led to $d_1$ as well. This
subsection shows that none does.

The standard construction reshapes the likelihood ratio of the belief against
the member of the model under which validity binds, here the worst member
$P^\star$, so
\[
  r=\frac{\widehat P_1}{P^\star}=(1.5188,\;0.9521,\;0.7433,\;0.7666).
\]
The ratio takes a different value at every outcome, so
Proposition~\ref{prop:no_shape} does not rule out writing the safe profile as
$w(r)$.

The obstacle is the second one described after that proposition. The
shape functions used in practice are increasing in $r$, so they rank the
outcomes as $r$ does,
\[
  a\ \prec\ a^{\mathrm e}\ \prec\ s^{\mathrm e}\ \prec\ s ,
\]
giving the most e-value to $s$, the outcome the belief rates likeliest
relative to the model. The loss of $d_1$ ranks them almost the other way,
\[
  s\ \prec\ s^{\mathrm e}\ \prec\ a\ \prec\ a^{\mathrm e} .
\]
To keep the certificate of $d_1$ low, the profile must put its largest
e-value on the two Argentina outcomes, where ordering mostly Spain is ruinous.
An increasing shape does the opposite. It puts its largest e-value on $s$,
where the order costs almost nothing, and its smallest on $a$ and
$a^{\mathrm e}$, where the order is most exposed. The safe profile
\eqref{eq:market_safe} reverses this, which is why it is not increasing in $r$.

Failing to match the safe profile exactly is not the same as failing to reach
$d_1$, so the question is settled by searching over shapes directly. Two
searches are run, both keeping every profile valid with its budget fully
spent, so that only the shape varies. The first covers the capped power
family \eqref{eq:power_shape} of Section~\ref{sec:result3}, sweeping $120$
values of $h$ in $[-6,0.95]$ against $60$ values of $\alpha$ in
$[0.01,0.95]$, a total of $7{,}200$ combinations, with $\lambda$ set by
bisection. The second drops the parametric form altogether. It draws
$200{,}000$ profiles by sampling four values uniformly, sorting them and
assigning them to the outcomes in increasing order of $r$, which gives a
random increasing shape.

Neither search ever reaches $d_1$. Every shape tried leads to the even split
or to mostly Argentina. Restricting to increasing shapes therefore does not
merely make $d_1$ more expensive to protect. It puts $d_1$ out of reach
altogether, while the safe profile reaches it with certificate $18$.

One input is left to the analyst, the historical rate at which finals go
beyond ninety minutes. Repeating the analysis for rates from $6/21$ to
$10/21$, two finals either side of the observed $8/21$, the two criteria
disagree at all five rates, and at all five no capped power shape reaches
$d_1$. The finding does not depend on that choice.

What the application cannot show is which criterion was right. Spain won
after extra time \citep{AlJazeera2026final}, the outcome $s^{\mathrm e}$, in which
ordering mostly Spain costs $3$ and the even split $19$. But a single match is
not a verdict on a decision rule. Settling that requires a setting where the
truth is known in advance.
